% Abstract (black-book template abstract.tex).
\clearpage
\thispagestyle{plain}
\addcontentsline{toc}{chapter}{Abstract}
\begin{center}
  {\Large\bfseries Abstract}
\end{center}
\vspace{0.6cm}

\noindent Idle computers are common in homes, offices and college laboratories,
but using them for a Python script normally requires a configured cluster, a
program rewritten against a cluster API, and machines that can reach one
another through their routers. This report presents \textbf{PipedPeer}, a
system that runs an unmodified Python script on other machines with none of
these. Each machine runs one daemon from a single binary; the submitting machine
builds the script's environment as a pinned \textbf{Nix closure}, places the job
by live capacity, sends only the store paths the chosen machine lacks, and runs
the job in a \texttt{crun} sandbox under its own memory ceiling. A shim injected
into the interpreter \textbf{intercepts the parallel primitives the script already
uses} --- \texttt{multiprocessing.Pool}, \texttt{joblib}, pandas, NumPy and torch
data-parallel training --- and spreads their work across machines in proportion
to measured throughput, only when a cost model predicts a gain. Machines on
different networks meet through an introducer that every daemon provides, then
connect directly over authenticated QUIC links after port mapping or hole
punching; no relay carries job data and no central server exists.

\noindent On a container laboratory, a job whose environment is already in place
starts in a median of 0.83\,s against 23.4\,s for its first submission. Between
two laptops in different homes, a mixed CPU script finished in 13.3--14.6\,s
against 17.4--18.1\,s on one laptop. On a two-rank training ring, synchronising
every few steps instead of every step kept the validation loss no more than 1\%
above one device's on every test model while sending 9 to 15 times fewer bytes. Of 829 tests,
822 pass and none fail, and a mutation audit that disables 95 features one at a
time is caught in every case.

\vspace{0.6cm}
\noindent\textbf{Keywords:} Distributed Computing, Peer-to-Peer, NAT Traversal,
Transparent Interception, Nix, Content-Addressed Caching, OCI Sandboxing,
Cost-Based Scheduling, Data-Parallel Training, Fault Tolerance.
